\documentclass[t]{beamer}
\usepackage{graphicx}
\usepackage[brazil]{babel}

\title{Análise de Desempenho: CFD em CUDA Fortran}
\author{Lucas Timm}
\date{31/08/2026}
\begin{document}

\maketitle

\newcommand{\inframe}[2]{
    \begin{frame}{#1}
        #2
    \end{frame}
}

\newcommand{\intwocols}[2]{
    \begin{columns}
        \begin{column}{0.5\textwidth}
            #1
        \end{column}
        \begin{column}{0.5\textwidth}
            #2
        \end{column}
    \end{columns}
}


\newcommand{\ttwocframe}[3]{
    \inframe{#1}{\intwocols{#2}{#3}}
}

\newcommand{\twocframe}[2]{
    \inframe{}{#1}{#2}
}

% Introdução CFD
\ttwocframe 
{Intro CFD}
{
    Computational Fluid Dynamics (CFD)  é o campo da mecânica computacional responsável por simular o comportamento de fluidos. Nesta campo de estudo as equações de Navier-Stokes, Equilíbrio de Massa e Momento, são discretizadas espacialmente por métodos numéricos como: Diferenças Finitas, Volumes Finitos, Elementos Finitos. Enquanto à discretização temporal são utilizados esquemas implícitos ou explícitos de marcha no tempo.
}
{
    As simulações de fluidos são conhecidas por necessitarem de malhas extremamente refinadas para a captura dos fenômenos com um grau de precisão adequado, visando um resultado útil. Para o trabalho atual será analisado um programa desenvolvido no PPGEC, Programa de Pós Graduação em Engenharia Civil, para aplicações na Engenharia do Vento Computacional (CWE). 
}

\ttwocframe
{Intro CFD}
{
    Por uma noção de ordem de grandeza, cada elemento hexaédrico corresponde a um sistema linear de 24 equações, podemos esperar que para um problema da vida real com 100.000 elementos, teríamos pelo menos 2.400.00 equações a serem resolvidas. E 100.000 elementos não é um número considerado muito grande, em um caso onde não há modelagem de turbulência (o que temos nesse caso), poderia ser até inviável resolver problemas reais com muita turbulência, devido a demanda computacional.
}
{
    Assim, as placas de vídeo se tornaram essenciais para melhorar o desempenho dos problemas de CWE, e CFD como um todo, devido ao declíneo relativo da capacidade de processamento de um core apenas, comparado com o ambiente massivamente paralelo da GPU. Neste trabalho será analisada a performance do programa, identificando o bottleneck específico, e aplicando alguma técnica de otimização.
}

\inframe
{Exemplos}
{
    \begin{figure}
        \centering
        \includegraphics[width = 1\textwidth]{./assets/pccfd.png}
        \caption{Simulação de CFD em um gabinete}
    \end{figure}
}

\inframe
{Exemplos}
{
    \begin{figure}
        \centering
        \centering
        \includegraphics[width = 1\textwidth]{./assets/planecfd.png}
        \caption{Simulação de CFD em um um avião}
    \end{figure}
}

\inframe
{Exemplos}
{
    \begin{figure}
        \centering
        \includegraphics[width = 1\textwidth]{./assets/cfdbuilding.jpg}
        \caption{Simulação de CFD de Estruturas}
    \end{figure}
}

\ttwocframe
{O Código}
{
     \begin{itemize}
        \item Desenvolvido por pós-graduandos da Engenharia Civil
        \item Elementos Finitos Hexaédricos Lineares
        \item Modelo Dinâmico de Turbulência
        \item CUDA Fortran
     \end{itemize}
}
{
     Está sendo utilizado atualmente no meu Trabalho de Conclusão de Curso na Engenharia Civil. 
     Em conversa com o doutorando que trabalhou no código, houve relatos que aumentar o número de thread por blocos não melhora a performance.
     Isso me faz pensar que o código pode ser memory-bound ou apenas não foi encontrado o número de threads otimizado.
}

\ttwocframe
{Metodologia}
{
    Criar uma malha de benchmark simples, idealmente com turbulência para ser um representativo de uma carga típica do código.
    Utilizar o ferramentário do CUDA Nvidia, o Nsight Compute, para analisar tempo de execução (por passo e total) kernel a kernel.

    Entender:
    \begin{itemize}
        \item Throughput
        \item Ocupância
        \item Tempo de Execução Total 
        \item Tempo de Execução por Passo de Tempo
    \end{itemize}

}
{
    Após identificado um gargalo, otimizar e fazer as medições novamente.

    \begin{figure}
        \centering
        \centering
        \includegraphics[width = 1\textwidth]{./assets/cavidade.png}
        \caption{Uma malha possível para análise.}
    \end{figure}
}

\ttwocframe{Metodologia}
{
A primeira etapa é provar ou desprovar a premissa que a performance é similar entre passos de tempo. Ou seja, se o bottleneck em um passo for o kernel 5, o kernel 5 sempre é o benchmark independente do passo de tempo analisado. Pode-se fazer isso registrando o tempo de execução de cada kernel por passo.

Identificado o kernel problemático, ele será extraído e testado com diversos números de threads diferentes para provar ou desprovar que isto afeta o desempenho.
}
{
Se for identificado algum outro problema no kernel, como os padrões de acesso à memória compartilhada e global da GPU, o livro, \textit{Programming Massively Parallel Processors}, traz uma tabela de checklist com algumas otimizações. Esta será a primeira fonte para as otimizações possíveis. O código será rodado no PCAD, assegurando o controle do ambiente de execução com o guix e slurm. Makefiles serão utilizados para garantir o workflow dos experimentos. 
}

\twocframe
{
    \begin{figure}
        \centering
        \includegraphics[width=0.75\textwidth]{./assets/pmpp.png}
        \caption{Principal referência para otimizações.}
    \end{figure}
}
{
    \begin{figure}
        \centering
        \centering
        \includegraphics[width=0.55\textwidth]{./assets/otimizacoes.png}
        \caption{Tabela de checklist de otimizações.}
    \end{figure}
}

\end{document}
